\documentclass[aspectratio=169,10pt]{beamer}
\usetheme{Madrid}
\usecolortheme{default}

\usepackage{hyperref}
\usepackage{tikz}
\usepackage{amsmath,amssymb}

% --- Custom Colors ---
\definecolor{unibblue}{RGB}{0,51,102}
\definecolor{unibgold}{RGB}{212,175,55}
\definecolor{unibgreen}{RGB}{0,128,0}

\setbeamercolor{palette primary}{bg=unibblue,fg=white}
\setbeamercolor{palette secondary}{bg=unibblue!85,fg=white}
\setbeamercolor{palette tertiary}{bg=unibblue!70,fg=white}
\setbeamercolor{structure}{fg=unibblue}
\setbeamercolor{block title}{bg=unibblue!85,fg=white}
\setbeamercolor{block body}{bg=unibblue!10,fg=black}
\setbeamercolor{alertblock title}{bg=unibgold!90,fg=black}
\setbeamercolor{alertblock body}{bg=unibgold!15,fg=black}

\title{Persamaan Diferensial}
\subtitle{Minggu 4: PDB Linier Orde 2 Homogen Koefisien Konstan}
\author{Novalio Daratha \& Muhammad Arfan}
\institute{Program Studi Teknik Elektro\\Universitas Bengkulu}
\date{2026}

\begin{document}

\begin{frame}
  \titlepage
\end{frame}

\begin{frame}{Capaian Pembelajaran (CPMK 2)}
  \begin{itemize}
    \item \textbf{CPMK 2.1:} Memahami bentuk standar PDB linear orde 2 homogen koefisien konstan.
    \item \textbf{CPMK 2.2:} Mampu menyusun dan menyelesaikan persamaan karakteristik / polinomial pembantu.
    \item \textbf{CPMK 2.3:} Menganalisis 3 jenis solusi basis: akar real berbeda, akar kembar, dan akar kompleks konjugat.
    \item \textbf{CPMK 2.4:} Mengaplikasikan prinsip superposisi dan menyelesaikan Masalah Nilai Awal (IVP) orde 2.
  \end{itemize}
\end{frame}

\begin{frame}{Daftar Isi}
  \tableofcontents
\end{frame}

\section{Bentuk Umum \& Persamaan Karakteristik}
\begin{frame}{Bentuk Standar PDB Orde 2 Homogen}
  \begin{block}{Persamaan Diferensial}
    $$ a \frac{d^2y}{dx^2} + b \frac{dy}{dx} + c y = 0, \quad a \neq 0 $$
    dengan $a, b, c$ adalah konstanta riil.
  \end{block}
  \pause
  \textbf{Strategi Solusi Eksponensial:}
  Asumsikan solusi berbentuk $y = e^{rx}$, maka $y' = r e^{rx}$ dan $y'' = r^2 e^{rx}$.
  \pause
  Substitusikan ke dalam PDB:
  $$ a(r^2 e^{rx}) + b(r e^{rx}) + c(e^{rx}) = 0 \implies e^{rx}(a r^2 + b r + c) = 0 $$
  Karena $e^{rx} \neq 0$, diperoleh:
  \begin{alertblock}{Persamaan Karakteristik}
    $$ {\color{unibblue} a r^2 + b r + c = 0} $$
    Akar-akar: $r_{1,2} = \frac{-b \pm \sqrt{b^2 - 4ac}}{2a}$.
  \end{alertblock}
\end{frame}

\section{Tiga Kasus Akar Karakteristik}
\begin{frame}{Kasus 1: Diskriminan $D > 0$ (Dua Akar Real Berbeda)}
  Jika $b^2 - 4ac > 0$, diperoleh dua akar real berbeda $r_1 \neq r_2$.
  \pause
  \begin{block}{Solusi Umum Kasus 1}
    $$ y(x) = C_1 e^{r_1 x} + C_2 e^{r_2 x} $$
  \end{block}
  \pause
  \textbf{Contoh Soal:} Selesaikan $y'' - 5y' + 6y = 0$.
  \begin{itemize}
    \item Persamaan karakteristik: $r^2 - 5r + 6 = 0 \implies (r-2)(r-3) = 0$.
    \item Akar-akar: $r_1 = 2, r_2 = 3$.
    \item Solusi umum: ${\color{unibblue} y(x) = C_1 e^{2x} + C_2 e^{3x}}$.
  \end{itemize}
\end{frame}

\begin{frame}{Kasus 2: Diskriminan $D = 0$ (Akar Real Kembar)}
  Jika $b^2 - 4ac = 0$, diperoleh akar kembar: $r_1 = r_2 = r = -\frac{b}{2a}$.
  \pause
  Hanya ada satu solusi bebas $y_1 = e^{rx}$. Menggunakan metode reduksi orde, solusi bebas kedua adalah $y_2 = x e^{rx}$.
  \pause
  \begin{block}{Solusi Umum Kasus 2}
    $$ y(x) = (C_1 + C_2 x)e^{rx} $$
  \end{block}
  \pause
  \textbf{Contoh Soal:} Selesaikan $y'' + 4y' + 4y = 0$.
  \begin{itemize}
    \item Karakteristik: $r^2 + 4r + 4 = (r+2)^2 = 0 \implies r = -2$.
    \item Solusi umum: ${\color{unibblue} y(x) = (C_1 + C_2 x)e^{-2x}}$.
  \end{itemize}
\end{frame}

\begin{frame}{Kasus 3: Diskriminan $D < 0$ (Akar Kompleks Konjugat)}
  Jika $b^2 - 4ac < 0$, diperoleh akar imajiner/kompleks:
  $$ r_{1,2} = \alpha \pm j\beta, \quad \alpha = -\frac{b}{2a}, \quad \beta = \frac{\sqrt{4ac - b^2}}{2a} $$
  \pause
  Berdasarkan Rumus Euler ($e^{j\theta} = \cos\theta + j\sin\theta$):
  \begin{block}{Solusi Umum Kasus 3 (Osilasi Teredam)}
    $$ y(x) = e^{\alpha x} \left( C_1 \cos(\beta x) + C_2 \sin(\beta x) \right) $$
  \end{block}
  \pause
  \textit{Keterangan Fisis: $\alpha$ menentukan laju redaman amplitudo, sedangkan $\beta$ adalah frekuensi osilasi sudut.}
\end{frame}

\section{Aplikasi Fisik: Rangkaian LC Murni}
\begin{frame}{Rangkaian LC Murni (Osilator Harmonik)}
  \begin{block}{Pemodelan KVL Rangkaian LC}
    $$ L \frac{d^2 i}{dt^2} + \frac{1}{C} i = 0 \implies \frac{d^2 i}{dt^2} + \omega_0^2 i = 0, \quad \omega_0 = \frac{1}{\sqrt{LC}} $$
  \end{block}
  \pause
  Persamaan karakteristik:
  $$ r^2 + \omega_0^2 = 0 \implies r = \pm j\omega_0 \quad (\alpha = 0, \beta = \omega_0) $$
  \pause
  \begin{alertblock}{Solusi Osilasi Tanpa Redaman}
    $$ i(t) = C_1 \cos(\omega_0 t) + C_2 \sin(\omega_0 t) $$
    Energi berpindah bolak-balik antara medan magnet induktor ($E_L = \frac{1}{2}Li^2$) dan medan listrik kapasitor ($E_C = \frac{1}{2}Cv^2$).
  \end{alertblock}
\end{frame}

\section{Kuis Interaktif}
\begin{frame}{Kuis Interaktif 2: Penyelesaian Masalah Nilai Awal (IVP)}
  \textbf{Soal:} Selesaikan IVP: $y'' + 9y = 0$ dengan syarat awal $y(0) = 4$ dan $y'(0) = 6$.
  \pause
  \begin{block}{Langkah Solusi:}
    \begin{enumerate}
      \item Karakteristik: $r^2 + 9 = 0 \implies r = \pm 3j$. \pause
      \item Solusi umum: $y(x) = C_1 \cos(3x) + C_2 \sin(3x)$. \pause
      \item Terapkan $y(0) = 4$: $C_1 \cos(0) + C_2 \sin(0) = 4 \implies C_1 = 4$. \pause
      \item Turunan: $y'(x) = -3C_1 \sin(3x) + 3C_2 \cos(3x)$. \pause
      \item Terapkan $y'(0) = 6$: $3C_2 = 6 \implies C_2 = 2$.
    \end{enumerate}
    $$ {\color{unibblue} y(x) = 4\cos(3x) + 2\sin(3x)} $$
  \end{block}
\end{frame}

\begin{frame}{Rangkuman Materi Minggu 4}
  \begin{table}[]
    \centering
    \begin{tabular}{|c|c|c|}
      \hline
      \textbf{Diskriminan} & \textbf{Akar Karakteristik} & \textbf{Bentuk Solusi Umum} \\ \hline
      $D > 0$ & Real berbeda: $r_1 \neq r_2$ & $y = C_1 e^{r_1 x} + C_2 e^{r_2 x}$ \\ \hline
      $D = 0$ & Real kembar: $r_1 = r_2 = r$ & $y = (C_1 + C_2 x)e^{rx}$ \\ \hline
      $D < 0$ & Kompleks: $\alpha \pm j\beta$ & $y = e^{\alpha x}(C_1 \cos\beta x + C_2 \sin\beta x)$ \\ \hline
    \end{tabular}
  \end{table}
\end{frame}

\end{document}
